\documentclass[CJKmath,a4paper,10pt]{ctexart}
\usepackage{geometry}
\geometry{inner=1.5cm,outer=1.5cm,top=2.45cm,bottom=1cm}
\usepackage[dvipsnames,svgnames,x11names,table]{xcolor}
\usepackage{exam-zh-choices}

\RequirePackage{amssymb} % Must be loaded before unicode-math
\RequirePackage{unicode-math} % Math fonts in xetexorluatex
%\setmathfont{texgyrepagella-math.otf}
%\setmainfont{STIXTwoText}[
%	Path=./fonts/STIXTwoText/,
%  Extension = .otf,
%  UprightFont = STIXTwoMath-Regular,
%  BoldFont = STIXTwoText-SemiBold,
%  ItalicFont = STIXTwoText-Italic,
%  BoldItalicFont = STIXTwoText-BoldItalic,
%]
%\setmathfont{STIXTwoMath-Regular.otf}

\usepackage{tikz,calc}
\usetikzlibrary{calc,arrows,shadows.blur,tikzmark}
% By default all math in TikZ nodes are set in inline mode. Change this to
% displaystyle so that we don't get small fractions.
\everymath{\displaystyle}

\usepackage[most]{tcolorbox}
\tcbuselibrary{breakable, skins,theorems}


\usepackage{xkeyval}
\makeatletter 
\usepackage{amsmath,amssymb,amsthm}
\usepackage{xkeyval} 

\tcbset{
    common/.style={
      fontupper=\rmfamily,
      lower separated=true,
      coltitle=black,
      colback=gray!5,
      boxrule=0.5pt,
      fonttitle=\bfseries,
      enhanced,
      breakable,
      top=8pt,
      before skip=8pt,
      attach boxed title to top left={yshift=-0.05in,xshift=0.15in},
      boxed title style={
      	boxrule=0pt,
        colframe=black,
        arc=0pt,
        outer arc=0pt,
        drop fuzzy shadow
     	},
      separator sign={.},
      drop fuzzy shadow
     },
     litistyle/.style={  %例题风格
     	common,
      colframe=gray,% 
      colback=white,
      colbacktitle=Gold,
      sharp corners,rounded corners=southeast,
      overlay unbroken and last={
      \node[anchor=south east, outer sep=0pt] at (\linewidth-width,0){\textcolor{black!60}{$\blacksquare$}};}
         },
% -----------------------------------------------------------------------------
% Explanation: BreakableBox@title key
% The following defines a pgfkeys/tcolorbox key named
%   "BreakableBox@title" which accepts 2 arguments (see 
%   "/.code n args={2}"). When the key is used like
%   "BreakableBox@title={<env>}{<optTitle>}" the code below is executed
%   with #1=<env> and #2=<optTitle>.
%
% Purpose:
% - If the second argument (#2) is empty, set the tcolorbox title to
%     \csname <env>name\endcsname~\thetcbcounter
%   which expands to the environment name (e.g. "liti" -> \litiname)
%   followed by the current counter value.
% - If the second argument is provided, append " (#2)" to the title.
%
% Notes:
% - This relies on \makeatletter being active so that keys with '@'
%   are allowed in control sequence names.
% - \ifblank is used to test whether #2 is empty; ensure the package
%   that provides \ifblank (for example, etoolbox or xparse) is loaded
%   earlier if necessary.
% - Example usage inside the code: \tcbset{title={...}} sets the
%   tcolorbox title key accordingly.
%
         BreakableBox@title/.code n args={2}
              {
                \ifblank{#2}
                  {\tcbset{title={\csname #1name\endcsname~\thetcbcounter}}}
                  {\tcbset{title={\csname #1name\endcsname~\thetcbcounter\ (#2)}}}
              },
}      
  % define an internal control sequence \newbreakablebox for fancy mode's newtheorem
  % #1 is the environment name, #2 is the prefix of label, #3 is the style
  % style: thmstyle, defstyle, prostyle
  % e.g. \newbreakablebox{theorem}{thm}{thmstyle}
  % will define two environments: numbered ``theorem'' and no-numbered ``theorem*''
  % WARNING FOR MULTILINGUAL: this cs will automatically find \theoremname's definition,
  % WARNING FOR MULTILINGUAL: it should be defined in language settings.
  \newcommand{\newbreakablebox}[3]{
    \ifcsundef{#1name}{%
      % define a default name if not defined before
      \tcbset{BreakableBox@title/.code n args={2}{\tcbset{title={use newcommand define #1name}}} }
    }{\relax}
    \DeclareTColorBox[auto counter,number within=section]{#1}{ g o t\label g }{
        common, % use common style
        #3, % use the style passed in
        IfValueTF={##1}
          {BreakableBox@title={#1}{##1}}
          {
            IfValueTF={##2}
            {BreakableBox@title={#1}{##2}}
            {BreakableBox@title={#1}{}}
          },
        IfValueT={##4}
          {
            IfBooleanTF={##3}
              {label={##4}}
              {VIVID@label={#2}{##4}}
          }
      }
    \DeclareTColorBox{#1*}{ g o }{
        common,#3,
        IfValueTF={##1}
          {BreakableBox@title={#1}{##1}}
          {
            IfValueTF={##2}
            {BreakableBox@title={#1}{##2}}
            {BreakableBox@title={#1}{}}
          },
      }
  }
  % define several environment 
  % we define headers like \definitionname before
  \newcommand{\litiname}{例题}
  \newbreakablebox{liti}{def}{litistyle}
  
  
%补充内容
%%设置新字体
%%定义带圈数字命令
\newfontfamily{\nmfont}{circlenumber}
[%
Extension=.otf,
Path=./fonts/]

\newcommand{\quan}[1]{{\nmfont \symbol{#1}}}
\newcommand{\kk}[1]{\quan{\numexpr32+#1}}%\kk{<参数范围1-95>}96、97、98、99分别用\quan{196} \quan{197} \quan{199} \quan{201}

%脚注使用带圈数字
\newcommand*\kkctr[1]{%
  \protect\kk{\number\numexpr\value{#1}\relax}}
\renewcommand*\thefootnote{\textcolor{black}{\kkctr{footnote}}}

%%无悬挂脚注格式
\renewcommand\@makefntext[1]{%
  \setlength\parindent{2\ccwd}\selectfont
  \@thefnmark\ #1}

%修改\part，使其不分页
\def\@endpart{%
	\thispagestyle{empty}
  \vskip40\p@%
   \@afterheading}



\RequirePackage{enumitem}
%\newenvironment{myenum}{\begin{enumerate}[label=\protect\kk{\arabic*}]\small}{\end{enumerate}}%
\setlist{noitemsep}
\setlist[enumerate, 1]{label=\protect\kk{\arabic*},itemsep=0.5ex}  
\makeatother



%%%marker环境
\newtcolorbox{marker}[1][]{enhanced,before skip=2mm,
	after skip=3mm,fontupper=\rmfamily,
	boxrule=0.4pt,left=5mm,right=2mm,top=1mm,bottom=1mm,
	colback=yellow!50,colframe=yellow!20!black,
	sharp corners,rounded corners=southeast,
	arc is angular,arc=3mm,underlay={%
		\path[fill=tcbcolback!80!black] ([yshift=3mm]interior.south east)--++(-0.4,-0.1)--++(0.1,-0.2);
		\path[draw=tcbcolframe,shorten <=-0.05mm,shorten >=-0.05mm] ([yshift=3mm]interior.south east)--++(-0.4,-0.1)--++(0.1,-0.2);
		\path[fill=yellow!50!black,draw=none] (interior.south west) rectangle node[white]{\Huge\bfseries !} ([xshift=4mm]interior.north west);
	},
	drop fuzzy shadow,#1
}
\begin{document}
\section{期中}
\begin{liti}
已知$f(x)$是$\mathbb{R}上单调函数，且$f$\big(f(x)-x^5-3x^3\big)=5$，则：
\begin{choices}
\item $f(-1)=-3$
\item $f(x)$是奇函数
\item $f(x)+f(-x)=2$
\item 不等式$f(x-1)>57$的解集是$(3,+\infty)$
\end{choices}
\tcblower
\textbf{解析：}

1. \textbf{确定函数 $f(x)$ 的解析式}

已知 $f(x)$ 是 $\mathbb{R}$ 上的单调函数，且满足方程：
\[ f\big(f(x)-x^5-3x^3\big)=5 \]

由于 $f(x)$ 是单调函数，对于方程 $f(t)=5$，如果 $f(x)$ 是严格单调的，那么 $t$ 只有唯一解。设这个解为 $u_0$，即 $f(u_0)=5$。

根据题意，括号内的部分必须恒等于 $u_0$：
\[ f(x) - x^5 - 3x^3 = u_0 \implies f(x) = x^5 + 3x^3 + u_0 \]

将 $f(x)$ 的表达式代入 $f(u_0)=5$ 中：
\[ f(u_0) = u_0^5 + 3u_0^3 + u_0 = 5 \]

令 $h(t) = t^5 + 3t^3 + t$，则 $h'(t) = 5t^4 + 9t^2 + 1 > 0$，说明 $h(t)$ 在 $\mathbb{R}$ 上单调递增。
观察可知，当 $t=1$ 时：
\[ 1^5 + 3(1)^3 + 1 = 1 + 3 + 1 = 5 \]
所以 $u_0 = 1$ 是唯一解。

因此，函数的解析式为：
\[ f(x) = x^5 + 3x^3 + 1 \]

验证单调性：$f'(x) = 5x^4 + 9x^2 \ge 0$，函数在 $\mathbb{R}$ 上单调递增，符合题意。

2. \textbf{逐一分析选项}

\begin{itemize}
    \item \textbf{A. $f(-1)=-3$}
    \[ f(-1) = (-1)^5 + 3(-1)^3 + 1 = -1 - 3 + 1 = -3 \]
    结论：正确

    \item \textbf{B. $f(x)$ 是奇函数}
    \[ f(-x) = (-x)^5 + 3(-x)^3 + 1 = -x^5 - 3x^3 + 1 \]
    \[ -f(x) = -x^5 - 3x^3 - 1 \]
    显然 $f(-x) \ne -f(x)$，且 $f(-x) \ne f(x)$。实际上 $f(x)-1$ 才是奇函数。
    结论：错误

    \item \textbf{C. $f(x)+f(-x)=2$}
    \[ f(x) + f(-x) = (x^5 + 3x^3 + 1) + (-x^5 - 3x^3 + 1) = 2 \]
    结论：正确

    \item \textbf{D. 不等式 $f(x-1)>57$ 的解集是 $(3,+\infty)$}
    首先寻找 $f(x)=57$ 的解。
    \[ x^5 + 3x^3 + 1 = 57 \implies x^5 + 3x^3 = 56 \]
    观察 $x=2$ 时：$2^5 + 3(2)^3 = 32 + 24 = 56$。
    所以 $f(2) = 57$。
    不等式化为 $f(x-1) > f(2)$。
    因为 $f(x)$ 是单调递增函数，所以：
    \[ x - 1 > 2 \implies x > 3 \]
    解集为 $(3, +\infty)$。
    结论：正确
\end{itemize}

综上所述，正确选项为 ACD。
\end{liti}

\begin{liti}
解题步骤
\begin{itemize}
\item   第一步：理解问题

函数定义为：$f(x) = \lg(x^2 - 2x + a)$ 其中对数函数 $\lg$ 的定义要求其参数为正数，即：$x^2 - 2x + a > 0$ 我们希望函数 \(f(x)\) 的值域为全体实数 $\mathbb{R}$。这意味着对于任意实数 y，都存在某个 x 使得：$\lg(x^2 - 2x + a) = y$

等价于：
$
x^2 - 2x + a = 10^y
$

由于 $10^y > 0$，问题转化为：二次函数
$g(x) = x^2 - 2x + a$
在满足 \(g(x) > 0\) 的 x 的取值范围内，其值域必须包含所有正实数，即 \((0, +\infty)\)。

\item  第二步：分析二次函数 \(g(x)\)

将 \(g(x)\) 完成平方：$g(x) = (x - 1)^2 + (a - 1)$

这是一个开口向上的抛物线，顶点在 x = 1 处，最小值为：$m = a - 1$

定义域为满足 \(g(x) > 0\) 的 x 的集合，即：

\[
D = \{ x \in \mathbb{R} \mid g(x) > 0 \}
\]

\item  第三步：确定 \(g(x)\) 在 D 上的值域

考虑三种情况：
\begin{enumerate}
\item  情况 1：a - 1 < 0，即 a < 1

此时最小值 m < 0，抛物线在 x = 1 处取负值，因此与 x 轴有两个交点，定义域 D 由两个不相交的区间组成。在每个区间上，\(g(x)\) 的值域都是 \((0, +\infty)\)，因此并集仍为 \((0, +\infty)\)。

$\checkmark$ 满足条件。

\item  情况 2：a - 1 = 0，即 a = 1

此时：$g(x) = (x - 1)^2$

最小值为 0，但定义域中排除 x = 1（因为 \(g(1) = 0\) 不满足 \(g(x) > 0\)）。

在 $x \ne 1$ 时，\(g(x)\) 的值域为 \((0, +\infty)\)。

$\checkmark$满足条件。

\item  情况 3：a - 1 > 0，即 a > 1

此时最小值 m > 0，抛物线完全在 x 轴上方，定义域为全体实数 $\mathbb{R}$。

但 \(g(x)\) 的值域为 \([m, +\infty)\)，不包含区间 \((0, m)\)，因此值域不包含所有正实数。

$\times$ 不满足条件。
\end{enumerate}
\item  第四步：总结

函数 \(f(x) = \lg(x^2 - 2x + a)\) 的值域为 $\mathbb{R}$ 当且仅当：$a \leq 1$

即：$a \in (-\infty, 1]$

\textbf{最终答案：}$\boxed{(-\infty,1]}$
\end{itemize}
\end{liti}
\end{document}